\providecommand{\weakconv}{\rightharpoonup}
\providecommand{\K}{\mathbb{K}}

\subsection{Изометричность вложения \(E\) в \(E^{**}\). Критерий слабой сходимости последовательности}

\begin{definition}[Слабая сходимость]
	Пусть \(E\) --- нормированное пространство. Последовательность
	\(\{x_n\}\subset E\) называется \textbf{слабо сходящейся} к \(x\in E\), если
	\[
	\forall f\in E^*
	\qquad
	f(x_n)\to f(x).
	\]
	Обозначение:
	\[
	x_n\weakconv x.
	\]
\end{definition}

\begin{definition}[Каноническое вложение \(E\) в \(E^{**}\)]
	Пусть \(E\) --- нормированное пространство. Для каждого \(x\in E\) определим функционал
	\[
	F_x:E^*\to\K,
	\qquad
	F_x(f)=f(x).
	\]
	Отображение
	\[
	J:E\to E^{**},
	\qquad
	Jx=F_x,
	\]
	называется \textbf{каноническим вложением} пространства \(E\) в \(E^{**}\).
\end{definition}

\begin{theorem}[Изометричность канонического вложения]
	Каноническое вложение
	\[
	J:E\to E^{**},
	\qquad
	(Jx)(f)=f(x),
	\]
	является линейной изометрией:
	\[
	\|Jx\|_{E^{**}}=\|x\|_E
	\qquad
	\forall x\in E.
	\]
\end{theorem}

\begin{proof}
	
	\textbf{1. Линейность отображения \(J\).}
	
	Для любых \(x,y\in E\), \(\alpha,\beta\in\K\) и \(f\in E^*\) имеем
	\[
	J(\alpha x+\beta y)(f)
	=
	f(\alpha x+\beta y)
	=
	\alpha f(x)+\beta f(y)
	=
	(\alpha Jx+\beta Jy)(f).
	\]
	Значит,
	\[
	J(\alpha x+\beta y)=\alpha Jx+\beta Jy.
	\]
	
	\textbf{2. Верхняя оценка нормы.}
	
	Для любого \(f\in E^*\) с \(\|f\|\leqslant 1\):
	\[
	|(Jx)(f)|
	=
	|f(x)|
	\leqslant
	\|f\|\,\|x\|
	\leqslant
	\|x\|.
	\]
	
	Следовательно,
	\[
	\|Jx\|
	=
	\sup_{\|f\|\leqslant 1}|(Jx)(f)|
	\leqslant
	\|x\|.
	\]о
	
	\textbf{3. Нижняя оценка нормы.}
	
	Если \(x=0\), то равенство очевидно. Пусть \(x\ne0\). По следствию из теоремы Хана--Банаха
	существует функционал \(f_0\in E^*\), такой что
	\[
	\|f_0\|=1,
	\qquad
	|f_0(x)|=\|x\|.
	\]
	Тогда
	\[
	\|Jx\|
	=
	\sup_{\|f\|\leqslant1}|(Jx)(f)|
	\geqslant
	|(Jx)(f_0)|
	=
	|f_0(x)|
	=
	\|x\|.
	\]
	
	Итак,
	\[
	\|Jx\|\leqslant \|x\|
	\qquad
	\text{и}
	\qquad
	\|Jx\|\geqslant \|x\|.
	\]
	Следовательно,
	\[
	\|Jx\|=\|x\|.
	\]
	Отображение \(J\) линейно и сохраняет норму, значит, \(J\) --- линейная изометрия.
\end{proof}

\begin{theorem}[Критерий слабой сходимости последовательности]
	Пусть \(E\) --- нормированное пространство.
	
	Последовательность \(x_n\) слабо сходится к \(x\) тогда и только
	тогда, когда выполнены два условия:
	
	\begin{enumerate}
		\item Последовательность норм ограничена:
		\(\exists M:\ \sup_n \|x_n\|\leqslant M.\)
		
		\item Сходимость есть на плотном множестве:
		\(\exists S\subset E^*,\ \overline{[S]}=E^*\) такое, что
		\(\forall s\in S:\ s(x_n)\to s(x).\)
	\end{enumerate}
\end{theorem}

\begin{proof}
	\textbf{1. Переход к операторам \(Jx_n\).}
	
	По определению слабой сходимости и канонического вложения \(J\):
	\[
	x_n\weakconv x
	\quad\Longleftrightarrow\quad
	\forall f\in E^*:\ 
	(Jx_n)(f)\to(Jx)(f).
	\]
	То есть
	\[
	x_n\weakconv x
	\quad\Longleftrightarrow\quad
	Jx_n\to Jx
	\quad\text{поточечно на }E^*.
	\]
	
	\textbf{2. Применение критерия поточечной сходимости.}
	
	Имеем
	\[
	Jx_n,Jx\in E^{**}=\mathcal{L}(E^*,\K).
	\]
	Пространство \(E^*\) банахово, а по условию
	\[
	\overline{[S]}=E^*.
	\]
	По критерию поточечной сходимости
	\[
	Jx_n\to Jx
	\]
	поточечно тогда и только тогда, когда
	\[
	\{\|Jx_n\|\}\ \text{ограничена}
	\]
	и
	\[
	(Jx_n)(s)\to(Jx)(s)
	\qquad
	\forall s\in S.
	\]
	
	\textbf{3. Возврат к \(x_n\).}
	
	Так как \(J\) --- изометрия и
	\[
	(Jx_n)(s)=s(x_n),
	\]
	получаем
	\[
	\{\|x_n\|\}\ \text{ограничена},
	\qquad
	s(x_n)\to s(x)
	\quad
	\forall s\in S.
	\]
	
	Следовательно,
	\[
	x_n\weakconv x
	\quad\Longleftrightarrow\quad
	\begin{cases}
		\{\|x_n\|\}\ \text{ограничена},\\
		s(x_n)\to s(x),\quad \forall s\in S.
	\end{cases}
	\]
\end{proof}